% Authored: architecture of the documentation system itself.
% validate-docs: allow-placeholder-words
% (this chapter names the placeholder markers the validator searches for)

\chapter{Documentation System}
\label{ch:documentation-system}

The manual you are reading is itself a build artifact of this repository. This
chapter documents how it is produced, because a documentation system that
cannot be regenerated on demand stops matching the code within weeks.

\section{Architecture}

% Authored TikZ: architecture of the documentation generator.

\begin{figure}[htbp]
\centering
\begin{tikzpicture}[
  src/.style={rectangle, draw=black, align=center, font=\scriptsize,
    text width=32mm, minimum height=7mm, inner sep=3pt},
  tool/.style={rectangle, draw=black, rounded corners=2pt, align=center,
    font=\scriptsize, text width=36mm, minimum height=7mm, inner sep=3pt,
    fill=shadegrey},
  model/.style={rectangle, draw=black, very thick, align=center,
    font=\scriptsize, text width=36mm, minimum height=7mm, inner sep=3pt},
  flow/.style={-{Latex[length=1.8mm]}, thick},
  weak/.style={-{Latex[length=1.8mm]}, rulegrey, dashed},
  node distance=3.6mm
]
  \node[src] (repo) {Repository source\\ \scriptsize Python, TOML, YAML, CSV, Make};

  \node[tool, below=8mm of repo] (inspect) {\texttt{inspect\_repository.py}\\ \scriptsize files, deps, CI, tests};
  \node[tool, left=6mm of inspect] (api) {\texttt{extract\_api.py}\\ \scriptsize AST only};
  \node[tool, right=6mm of inspect] (deps) {\texttt{extract\_dependencies.py}\\ \scriptsize import graph};

  \node[model, below=8mm of inspect] (jsonmodel) {Normalized documentation model\\ \scriptsize \texttt{docs/metadata/*.json}};

  \node[tool, below=8mm of jsonmodel] (rarch) {\texttt{render\_architecture.py}};
  \node[tool, left=6mm of rarch] (rref) {\texttt{render\_reference.py}};
  \node[tool, right=6mm of rarch] (reng) {\texttt{render\_engineering.py}};

  \node[model, below=8mm of rarch] (frags) {Generated LaTeX fragments\\ \scriptsize \texttt{docs/latex/generated/**}};
  \node[src, below=7mm of frags] (mainrepo) {\texttt{docs/latex/main.tex}\\ \scriptsize + authored chapters};
  \node[src, below=7mm of mainrepo] (pdf) {PDF\\ \scriptsize built separately by latexmk};

  \node[tool, right=10mm of frags] (validate) {\texttt{validate\_docs.py}\\ \scriptsize labels, refs, escaping,\\ coverage, secrets, manifest};
  \node[tool, left=10mm of frags] (latexutil) {\texttt{latex\_utils.py}\\ \scriptsize escaping + table primitives};

  \draw[flow] (repo) -- (inspect);
  \draw[flow] (repo.west) to[out=180, in=90] (api.north);
  \draw[flow] (repo.east) to[out=0, in=90] (deps.north);
  \draw[flow] (api) -- (jsonmodel);
  \draw[flow] (inspect) -- (jsonmodel);
  \draw[flow] (deps) -- (jsonmodel);
  \draw[flow] (jsonmodel) -- (rref);
  \draw[flow] (jsonmodel) -- (rarch);
  \draw[flow] (jsonmodel) -- (reng);
  \draw[flow] (rref) -- (frags);
  \draw[flow] (rarch) -- (frags);
  \draw[flow] (reng) -- (frags);
  \draw[flow] (frags) -- (mainrepo);
  \draw[flow] (mainrepo) -- (pdf);
  \draw[weak] (latexutil) -- (frags);
  \draw[weak] (frags) -- (validate);
  \draw[weak] (jsonmodel.east) to[out=0, in=90] (validate.north);
\end{tikzpicture}
\caption{The documentation pipeline. Shaded boxes are tools, heavy boxes are
intermediate artifacts. The JSON model is the only interface between the
extraction and rendering layers: nothing below it reads Python source, and
nothing above it emits markup.}
\label{fig:documentation-pipeline}
\end{figure}


Extraction and rendering are separate layers and never mix. Extractors read
source and configuration and emit JSON; renderers read that JSON and emit
LaTeX. A renderer never parses Python and an extractor never emits markup,
which means an incorrect signature is an extractor bug and a broken table is a
renderer bug --- the two are independently testable, and both are tested.

\begin{description}
  \item[\srcfile{docs/tools/inspect\_repository.py}] Files, languages,
    directory roles, dependencies, tool configuration, Make targets, CI
    workflows, containers, deployments, tests, datasets, the environment
    template, and quality gates $\rightarrow$
    \srcfile{repository-inventory.json}.
  \item[\srcfile{docs/tools/extract\_api.py}] Modules, classes, functions,
    methods, attributes, constants, type aliases, CLI commands, HTTP routes,
    and settings fields $\rightarrow$ \srcfile{api-index.json}.
  \item[\srcfile{docs/tools/extract\_dependencies.py}] The internal import
    graph with per-edge context, layer assignment, coupling metrics, and cycle
    analysis $\rightarrow$ \srcfile{module-dependencies.json}.
  \item[\srcfile{docs/tools/latex\_utils.py}] Pure rendering primitives:
    escaping, labels, long tables, description lists, listings, index entries.
  \item[\srcfile{docs/tools/render\_reference.py},
    \srcfile{render\_architecture.py}, \srcfile{render\_engineering.py}]
    Model $\rightarrow$ LaTeX fragments under
    \srcfile{docs/latex/generated/}.
  \item[\srcfile{docs/tools/generate\_docs.py}] Orchestration and the
    manifest.
  \item[\srcfile{docs/tools/validate\_docs.py}] The checks in
    \cref{sec:doc-validation}.
\end{description}

\section{Safety properties}

\begin{description}
  \item[No application import.] Extraction uses \texttt{ast} only. Generating
    documentation cannot open a network connection, contact an LLM provider,
    load a model, build an index, or write outside the documentation tree. The
    only subprocess invoked is \texttt{git}, for revision facts.
  \item[Deterministic output.] Objects are emitted in fully-qualified-name
    order; timestamps come from the HEAD commit date rather than the wall
    clock. Two runs over one revision are byte-identical, which is what makes
    the manifest digests meaningful.
  \item[Total escaping.] Every string that reaches LaTeX passes through
    \texttt{escape\_latex}, which handles the ten special characters plus
    angle brackets and the vertical bar. Docstrings and source text cannot
    break the build.
  \item[Generated files are disposable.] \texttt{make docs} deletes and
    rewrites everything under \srcfile{docs/latex/generated/}, so an object
    removed from the source disappears from the manual rather than lingering.
    Authored files are never touched.
\end{description}

\section{Validation}
\label{sec:doc-validation}

\texttt{make docs} runs the validator automatically. It checks:

\begin{enumerate}[nosep]
  \item every public class and function in the model has a rendered fragment;
  \item every rendered signature matches the extracted signature;
  \item every LaTeX label is unique across the whole tree;
  \item every \texttt{\textbackslash cref} and \texttt{\textbackslash ref}
    target is defined somewhere;
  \item every \texttt{\textbackslash input} target exists on disk;
  \item generated files contain no unescaped LaTeX special characters outside
    verbatim contexts;
  \item every required authored file exists;
  \item no placeholder text (\texttt{TODO}, \texttt{FIXME}, \texttt{TBD},
    \texttt{lorem ipsum}, \texttt{XXX}) remains;
  \item no credential-shaped literal appears in any generated file;
  \item the manifest digests match the files on disk, and the manifest revision
    matches the current \texttt{HEAD}.
\end{enumerate}

Failures are reported per check with a count and a sample, and the exit status
is non-zero, so the validator is usable as a CI gate.

\section{Regeneration and CI}

\begin{lstlisting}[language=bash]
make docs        # regenerate and validate
make docs-check  # validate only, no regeneration
make docs-pdf    # regenerate, validate, then latexmk -pdf
\end{lstlisting}

The CI workflow gained a \texttt{documentation} job that runs
\texttt{make docs} and then fails if regeneration changed any tracked file.
That single check enforces the invariant that matters: the committed
documentation corresponds to the committed source. It does not compile the PDF,
because no LaTeX toolchain is installed in CI and the canonical deliverable is
the LaTeX source.

\section{Tests}

\srcfile{tests/test\_docs\_tools.py} covers the documentation system with the
same discipline as the rest of the repository: escaping of every special
character, idempotence properties of the escaper, table rendering, label
sanitisation, docstring section parsing, signature rendering including
keyword-only markers and defaults, exception collection, CLI and endpoint
extraction against the real source, cycle detection on synthetic graphs, and an
end-to-end assertion that the extracted model covers the actual public surface.
